\documentclass[12pt]{article}
\usepackage[margin=1in]{geometry}
\usepackage{enumitem}
\usepackage{hyperref}

\title{Using the \texttt{LLM-DRL restructured code} Project}
\author{}
\date{}

\begin{document}
\maketitle

\section{Overview of the Code Base}
\begin{itemize}[leftmargin=1.2cm]
  \item The project converts the original notebook pipeline into regular Python modules so that training and evaluation can be run from the command line.
  \item All source files live inside the directory \texttt{LLM-DRL restructured code}. The root project (one level above) still stores data, logs, and memory files just as before.
  \item Configuration values (paths, hyper-parameters, overwrite behaviour) are centralised in \texttt{configs.py}. Both \texttt{train.py} and \texttt{evaluate.py} read from this file so that there is a single place to adjust settings.
  \item The package \texttt{rl\_llm} contains the environment, PPO implementation, LLM controller, and supporting utilities copied from the notebook with only minimal glue changes.
\end{itemize}

\section{Roles of the Individual Python Files}
\subsection{Top-Level Files}
\begin{itemize}[leftmargin=1.2cm]
  \item \texttt{configs.py}: defines constants such as run names, log directories, reward weights, schedule parameters, memory paths, and helper \texttt{ensure\_unique\_subdir}. Edit this file to change the behaviour for new runs.
  \item \texttt{train.py}: the main training script. It creates the environment, PPO agent, LLM controller, selects (or creates) a log directory, resumes from checkpoints when present, and runs the PPO loop. It is the entry point for training.
  \item \texttt{evaluate.py}: the standalone evaluator. It loads a saved PPO policy, runs deterministic rollouts, captures frames, and saves GIFs. It never writes into the training log directory.
  \item \texttt{README.md}: short human-readable summary of the restructure.
  \item \texttt{LLM\_DRL\_documentation.tex} (this document): full instructions in \LaTeX{} format.
\end{itemize}

\subsection{Package \texttt{rl\_llm}}
\begin{itemize}[leftmargin=1.2cm]
  \item \texttt{env.py}: builds the scenario data from \texttt{featurefile\_transformed\_2.csv}, defines the \texttt{Agent} class (ownship \& intruder) and the Gym \texttt{StructuredEnv}, handles the reward shaping, and implements rendering.
  \item \texttt{ppo.py}: contains running-statistics utilities, rollout buffer, hybrid actor-critic network with attention, PPO algorithm class, and checkpoint helpers. The training script instantiates a \texttt{PPO} object from here.
  \item \texttt{llm.py}: stores prompt building utilities, the LLM controller finite state machine, ghost comparison logic, memory management helpers, and wrapper functions that call the local LLM. \texttt{train.py} relies on this module during LLM-assisted episodes.
  \item \texttt{utils.py}: holds geometry helpers, sector plotting code, and scenario selection utilities previously defined in the notebook.
  \item \texttt{\_\_init\_\_.py}: exports convenience symbols so modules can be imported more easily.
\end{itemize}

\section{Running \texttt{train.py}}
\subsection{What Happens During Training}
\begin{enumerate}[leftmargin=1.2cm]
  \item \textbf{Log Directory Creation.} The script calls \texttt{ensure\_unique\_subdir(LOG\_ROOT, RUN\_NAME)}. If a directory with the same run name exists, a timestamp suffix is added automatically (for example \texttt{log/PPO/PPO\_v7\_LLM\_RL\_20240311T101500}). This behaviour prevents accidental overwrites unless \texttt{ALLOW\_TRAIN\_OVERWRITE} is set to \texttt{True}.
  \item \textbf{Checkpoint Files.} Inside the chosen log directory the following files appear over time:
    \begin{itemize}
      \item \texttt{best\_model.pth}: best-performing policy weights (no optimizer state).
      \item \texttt{last.ckpt}: full checkpoint with optimizer, observation normaliser, and RNG states for resuming.
      \item \texttt{checkpoint\_\{step\}.pth}: periodic weight-only snapshots created every \texttt{SAVE\_MODEL\_FREQ} steps.
      \item \texttt{events.out.tfevents.*}: TensorBoard event files generated by \texttt{SummaryWriter}.
    \end{itemize}
  \item \textbf{LLM Memory Folders.} Whenever an episode uses the LLM controller, a fresh directory \texttt{memory/run\_YYYYMMDDTHHMMSS\_\_epXXXXXX} is created. It keeps the JSON prompt, decision, and ghost-comparison data. The helper \texttt{prune\_memory\_runs} retains the newest fifteen and deletes older ones.
  \item \textbf{Console Output.} The script prints the log folder path, whether it resumed from a checkpoint, evaluation summaries, and memory pruning events. This makes it easy to watch progress without opening files.
\end{enumerate}

\subsection{Changing Training Parameters}
\begin{itemize}[leftmargin=1.2cm]
  \item \textbf{Hyper-parameters.} Edit \texttt{configs.py}: the PPO schedule (\texttt{MAX\_TIMESTEPS}, \texttt{UPDATE\_TIMESTEP}, \texttt{EVAL\_FREQ}, \texttt{SAVE\_MODEL\_FREQ}, etc.), environment reward weights (\texttt{REWARD\_PARAMS}), LLM usage pattern (\texttt{LLM\_EVERY\_K}, \texttt{LLM\_FIRST\_N}), and the run name (\texttt{RUN\_NAME}).
  \item \textbf{Starting From Scratch.} Because the log directory name is unique, a new run already starts fresh. To reuse an existing name, either delete the old log folder or set \texttt{ALLOW\_TRAIN\_OVERWRITE = True} (only do this when you intentionally want to continue inside the same directory). If you want to ignore existing checkpoints while keeping past logs, change \texttt{RUN\_NAME} to a new value or delete \texttt{last.ckpt} in the chosen directory before starting.
  \item \textbf{Changing Memory Retention.} Adjust the \texttt{keep} value in the call \texttt{prune\_memory\_runs(base\_dir=MEMORY\_DIR, keep=\ldots)} inside \texttt{train.py}.
  \item \textbf{Dataset Path.} If the feature file moves, update \texttt{FEATUREFILE\_PATH} in \texttt{configs.py}.
\end{itemize}

\section{Running \texttt{evaluate.py}}
\subsection{What Happens During Evaluation}
\begin{enumerate}[leftmargin=1.2cm]
  \item \textbf{Selecting the Policy.} The script constructs \texttt{best\_model\_path = LOG\_ROOT / EVAL\_RUN\_NAME / ``best\_model.pth''}. Make sure \texttt{EVAL\_RUN\_NAME} matches the folder produced by the training run you want to inspect. If the file is missing, the script aborts with an error.
  \item \textbf{Output Directory.} GIFs are written to a unique subdirectory of \texttt{EVAL\_OUTPUT\_ROOT}. The helper \texttt{ensure\_unique\_subdir} again prevents overwriting (unless \texttt{ALLOW\_EVAL\_OVERWRITE} is set to \texttt{True}). Each evaluation invocation creates a folder such as \texttt{Evaluation GiFs/PPO\_v4\_eval\_20240311T104200}.
  \item \textbf{Intermediate Frames.} For every episode, PNG frames are saved in a sibling directory (e.g. \texttt{ep001/image\_000.png}) before being combined into a GIF. Old frames inside the same folder are cleared at the start of each episode.
  \item \textbf{Console Output.} After each evaluation episode the script prints the GIF location and return. A final line reports the mean return across all sampled episodes.
\end{enumerate}

\subsection{Changing Evaluation Parameters}
\begin{itemize}[leftmargin=1.2cm]
  \item \textbf{Policy Folder.} Edit \texttt{EVAL\_RUN\_NAME} in \texttt{configs.py} so that it points to the desired training run directory.
  \item \textbf{Overwrite Behaviour.} Set \texttt{ALLOW\_EVAL\_OVERWRITE = True} if you prefer to reuse an existing output folder instead of generating a timestamped one.
  \item \textbf{Number of Episodes.} Modify the loop at the bottom of \texttt{evaluate.py} (the \texttt{for ep in range(1, 2):} line) to save more or fewer GIFs.
  \item \textbf{GIF Frame Rate.} Adjust the \texttt{fps} argument in \texttt{\_save\_gif}.
\end{itemize}

\section{Suggested Workflow}
\begin{enumerate}[leftmargin=1.2cm]
  \item Review \texttt{configs.py} and set run-specific values (run name, hyper-parameters, dataset location).
  \item Launch \texttt{python train.py}. Watch the console for the chosen log directory and periodic evaluation summaries. Use \texttt{tensorboard --logdir log/PPO} to visualise metrics.
  \item After training, note the run directory (for example \texttt{log/PPO/PPO\_v7\_LLM\_RL\_20240311T101500}) and set \texttt{EVAL\_RUN\_NAME} to match it.
  \item Run \texttt{python evaluate.py} to build GIFs. Inspect the generated folder under \texttt{Evaluation GiFs}.
\end{enumerate}

\section{Final Remarks}
\begin{itemize}[leftmargin=1.2cm]
  \item The restructure keeps every component in plain Python files, so the environment can be imported into other projects or unit tests if needed.
  \item All file paths resolve relative to the project root, ensuring that data and logs integrate with the existing directory structure.
  \item Prompt strings for the LLM controller must still be supplied manually in \texttt{rl\_llm/llm.py}. Once inserted, the controller records each decision inside the \texttt{memory/} folder hierarchy.
\end{itemize}

\end{document}
